% =====================================================================
% Fiche de revision - Python Scripting - Lesson 5 : FastAPI
% Compilation : pdfLaTeX (Overleaf), aucune option particuliere requise.
% =====================================================================

\documentclass[a4paper,10pt]{article}

\usepackage[utf8]{inputenc}
\usepackage[T1]{fontenc}
\usepackage[french]{babel}
\usepackage[a4paper,margin=1.3cm]{geometry}
\usepackage{xcolor}
\usepackage{multicol}
\usepackage{booktabs}
\usepackage{array}
\usepackage{tabularx}
\usepackage{enumitem}
\usepackage{listings}
\usepackage{titlesec}
\usepackage{tcolorbox}

% ---------------------------------------------------------------- couleurs
\definecolor{accent}{HTML}{1F5FA8}
\definecolor{codefg}{HTML}{A03030}
\definecolor{codebg}{HTML}{F4F6F8}
\definecolor{kw}{HTML}{0B5FA5}
\definecolor{cmt}{HTML}{6C8A54}
\definecolor{strcol}{HTML}{2E7D32}
\definecolor{danger}{HTML}{B3261E}

% ---------------------------------------------------------------- titres
\titleformat{\section}
  {\normalfont\large\bfseries\color{accent}}{}{0pt}{}
  [\vspace{-0.75em}{\color{accent}\rule{\linewidth}{0.7pt}}]
\titlespacing*{\section}{0pt}{1.4ex}{0.7ex}

\titleformat{\subsection}[runin]
  {\normalfont\bfseries}{}{0pt}{}[\;---\;]
\titlespacing*{\subsection}{0pt}{0.8ex}{0.4em}

% ---------------------------------------------------------------- code
\lstset{
  language=Python,
  basicstyle=\ttfamily\footnotesize,
  keywordstyle=\color{kw}\bfseries,
  commentstyle=\color{cmt}\itshape,
  stringstyle=\color{strcol},
  showstringspaces=false,
  columns=fullflexible,
  keepspaces=true,
  backgroundcolor=\color{codebg},
  frame=none,
  xleftmargin=0.3em,
  aboveskip=0.55em,
  belowskip=0.55em,
  breaklines=true,
}

% code en ligne
\newcommand{\co}[1]{\texttt{\textcolor{codefg}{#1}}}

% ---------------------------------------------------------------- boites
\newtcolorbox{piegebox}[1]{
  colback=red!3, colframe=danger, boxrule=0.7pt,
  left=5pt, right=5pt, top=4pt, bottom=4pt,
  title=#1, fonttitle=\bfseries\small, coltitle=white
}

\newtcolorbox{astucebox}[1]{
  colback=accent!5, colframe=accent, boxrule=0.7pt,
  left=5pt, right=5pt, top=4pt, bottom=4pt,
  title=#1, fonttitle=\bfseries\small, coltitle=white
}

\setlist[itemize]{leftmargin=1.1em, itemsep=0.15em, topsep=0.25em, parsep=0pt}
\setlength{\columnsep}{1.1em}
\setlength{\parindent}{0pt}
\renewcommand{\arraystretch}{1.15}


\begin{document}

% Desactive les espaces automatiques du francais avant : ; ! ?
% pour que le code en ligne (ex. "item_id: int") ne soit pas deforme.
\shorthandoff{;:!?}

\small

\begin{center}
  {\LARGE\bfseries\color{accent} Python Scripting --- Lesson 5}\\[0.3em]
  {\normalsize FastAPI, SQLAlchemy, Pydantic \quad$\cdot$\quad fiche de révision QCM}
\end{center}
\vspace{0.4em}

% =====================================================================
\begin{piegebox}{Les 16 réponses pièges à connaître par c\oe ur}
\begin{minipage}[t]{0.485\linewidth}
\begin{tabularx}{\linewidth}{@{}l>{\raggedright\arraybackslash}X@{}}
\toprule
\textbf{Situation} & \textbf{Résultat} \\
\midrule
\co{main:app}                 & \textbf{module} \co{main}, \textbf{variable} \co{app} \\
\co{/items/42}, sans type     & \co{"42"} : une \textbf{chaîne} \\
\co{item\_id: int} et \co{abc}& \textbf{422}, la fonction n'est pas appelée \\
champ JSON manquant           & \textbf{422} aussi (pas 400) \\
\co{/users/\{id\}} avant \co{/users/me} & \co{me} part dans \co{id} : 422 \\
POST sans \co{status\_code}   & \textbf{200}, pas 201 \\
bonne URL, mauvaise méthode   & \textbf{405} (URL inconnue : 404) \\
\co{.first()} sans résultat   & \co{None} $\to$ \co{null} avec \textbf{200} \\
\bottomrule
\end{tabularx}
\end{minipage}\hfill
\begin{minipage}[t]{0.485\linewidth}
\begin{tabularx}{\linewidth}{@{}l>{\raggedright\arraybackslash}X@{}}
\toprule
\textbf{Situation} & \textbf{Résultat} \\
\midrule
\co{return HTTPException(...)} & renvoyé en 200 : il faut \co{raise} \\
\co{add()} sans \co{commit()}  & rien en base, \co{id} vaut \co{None} \\
email en double (\co{unique}) & \co{IntegrityError} $\to$ \textbf{500} \\
\co{Optional[str]} sans \co{= None} & \textbf{requis} (mais accepte \co{None}) \\
\co{response\_model} sans \co{id} & l'\co{id} est \textbf{retiré} de la réponse \\
\co{create\_all} + nouvelle colonne & table existante \textbf{inchangée} \\
\co{Depends(get\_db())}       & \co{TypeError} : passer la \textbf{fonction} \\
\co{Column(String)}           & \textbf{nullable} par défaut \\
\bottomrule
\end{tabularx}
\end{minipage}
\end{piegebox}

\vspace{0.5em}
\begin{multicols}{2}

% =====================================================================
\section{Installation et lancement}

\begin{lstlisting}
pip install fastapi uvicorn sqlalchemy
uvicorn main:app --reload     # depuis le dossier de main.py
\end{lstlisting}

\begin{itemize}
  \item \textbf{FastAPI} : le framework. \textbf{Uvicorn} : le serveur
        \textbf{ASGI} qui l'exécute (ASGI = asynchrone, WSGI = synchrone).
  \item Adresse par défaut : \co{http://127.0.0.1:8000}.
  \item \co{--reload} : redémarre à chaque modification ; \textbf{développement}
        uniquement.
  \item \co{main:app} est un import : lancé depuis un autre dossier,
        \co{Could not import module "main"}.
\end{itemize}

\subsection{Documentation automatique}
\co{/docs} (Swagger UI, bouton \emph{Try it out}), \co{/redoc} (ReDoc), et
le schéma brut \co{/openapi.json}. Tout est déduit des décorateurs, des
annotations de type et des schémas Pydantic : rien à écrire à la main.

% =====================================================================
\section{Première route}

\begin{lstlisting}
from fastapi import FastAPI

app = FastAPI()               # l'application

@app.get("/")                 # GET sur /
async def home():
    return {"message": "Hello World"}   # dict -> JSON
\end{lstlisting}

\subsection{D'où vient chaque paramètre ?} FastAPI regarde le
\textbf{type}, pas le nom :

\begin{tabularx}{\linewidth}{@{}l>{\raggedright\arraybackslash}X@{}}
\toprule
nom entre \co{\{\}} dans le chemin & \textbf{path} \\
type simple hors du chemin         & \textbf{query} (\co{?limit=5}) \\
modèle Pydantic                    & \textbf{body} JSON \\
annoté \co{Request}                & l'objet requête \\
\co{Depends(f)}                    & résultat de \co{f} \\
\bottomrule
\end{tabularx}

\vspace{0.3em}
\begin{lstlisting}
@app.get("/items/{item_id}")
async def read_item(item_id):        # "42" (str)
@app.get("/items_type/{item_id}")
async def read_item_typed(item_id: int):   # 42
@app.get("/items")
def list_items(limit: int = 10):     # ?limit=5, optionnel
\end{lstlisting}

\subsection{Ordre} les routes sont testées dans l'ordre de déclaration :
les chemins fixes (\co{/users/me}) \textbf{avant} les chemins à paramètre.

\subsection{\co{async def} ou \co{def}} les deux marchent. Les \co{def}
tournent dans un thread à part : à préférer quand le code est bloquant
(SQLAlchemy synchrone).

% =====================================================================
\section{REST et codes HTTP}

\begin{tabularx}{\linewidth}{@{}lll>{\raggedright\arraybackslash}X@{}}
\toprule
\textbf{CRUD} & \textbf{Méthode} & \textbf{Cours} \\
\midrule
Create & POST   & \co{/adduser} \\
Read   & GET    & \co{/users}, \co{/user/\{user\_name\}} \\
Update & PUT    & \co{/user/\{user\_id\}} \\
Delete & DELETE & \co{/user/\{user\_id\}} \\
\bottomrule
\end{tabularx}

\vspace{0.3em}
\textbf{200} OK \quad \textbf{201} Created \quad \textbf{404} Not Found
\quad \textbf{405} Method Not Allowed \quad \textbf{409} Conflict \quad
\textbf{422} validation \quad \textbf{500} bug serveur.

% =====================================================================
\section{SQLAlchemy : \co{database.py}}

\begin{lstlisting}
from sqlalchemy import create_engine
from sqlalchemy.orm import declarative_base, sessionmaker

SQLALCHEMY_DATABASE_URL = "sqlite:///./test.db"
engine = create_engine(SQLALCHEMY_DATABASE_URL)
SessionLocal = sessionmaker(autocommit=False,
                            autoflush=False, bind=engine)
Base = declarative_base()
\end{lstlisting}

\begin{itemize}
  \item \textbf{ORM} : une classe = une table, un objet = une ligne.
  \item \co{engine} : connexions ; \co{SessionLocal()} : \textbf{une
        nouvelle} session ; \co{Base} : classe mère des modèles.
  \item \co{./test.db} : relatif au dossier \textbf{courant}.
  \item \co{autocommit=False} : rien n'est écrit sans \co{commit()}.
  \item Cours : \co{sqlalchemy.ext.declarative} $\to$ avertissement en 2.x,
        importer depuis \co{sqlalchemy.orm}.
\end{itemize}

% =====================================================================
\section{Modèle : \co{models.py}}

\begin{lstlisting}
class User(Base):
    __tablename__ = "users"          # nom de la TABLE
    id = Column(Integer, primary_key=True, index=True)
    name = Column(String)
    email = Column(String, unique=True)
    nickname = Column(String)
    url = Column(String, nullable=True)   # deja le defaut

Base.metadata.create_all(bind=engine)    # dans main.py
\end{lstlisting}

\begin{itemize}
  \item \co{primary\_key} Integer : rempli automatiquement (1, 2, 3\dots).
  \item \co{unique=True} : vérifié par la \textbf{base}, au \co{commit}.
  \item \co{create\_all} crée les tables \textbf{manquantes}, ne modifie
        jamais une table existante (pas de migration : supprimer
        \co{test.db} ou utiliser Alembic).
\end{itemize}

% =====================================================================
\section{Pydantic : \co{schemas.py}}

\begin{lstlisting}
class UserSchema(BaseModel):
    name: str                        # requis
    email: str
    nickname: str
    url: Optional[str] = None        # optionnel (defaut)
    model_config = ConfigDict(from_attributes=True)

class UserOut(UserSchema):           # sortie, avec l'id
    id: int
\end{lstlisting}

\begin{itemize}
  \item Valide le \textbf{corps} reçu et décrit la \textbf{réponse}. Le
        modèle SQLAlchemy, lui, décrit la \textbf{table}.
  \item \co{orm\_mode = True} (Pydantic 1) = \co{from\_attributes=True}
        (Pydantic 2) : lire \co{user.name} sur un objet, pas un dict.
  \item Optionnel = a une \textbf{valeur par défaut}. \co{Optional[str]}
        veut seulement dire \og \co{str} ou \co{None} \fg.
  \item Conversion : \co{"42"} et \co{4.0} $\to$ \co{int} OK ;
        \co{"4.5"}, \co{"abc"} refusés ; \co{123} pour un \co{str} refusé.
\end{itemize}

% =====================================================================
\section{Session par requête : \co{Depends}}

\begin{lstlisting}
def get_db():
    db = SessionLocal()
    try:
        yield db          # donnee a la route
    finally:
        db.close()        # toujours, meme sur exception

@app.post("/adduser", response_model=UserOut)
def add_user(url: Request, request: UserSchema,
             db: Session = Depends(get_db)):
    user = User(name=request.name, email=request.email,
                nickname=request.nickname,
                url=f"{url.base_url}user/{request.name}")
    db.add(user)          # dans la session
    db.commit()           # INSERT
    db.refresh(user)      # relit : user.id
    return user
\end{lstlisting}

Dans le cours, \co{url} est la \co{Request} et \co{request} est le corps :
seuls les \textbf{types} comptent. \co{base\_url} vaut
\co{http://127.0.0.1:8000/}.

% =====================================================================
\section{Lire, modifier, supprimer}

\begin{tabularx}{\linewidth}{@{}l>{\raggedright\arraybackslash}X@{}}
\toprule
\co{db.query(User).all()} & une \textbf{liste} (vide si rien) \\
\co{.filter(User.name == n)} & condition SQL : \co{==}, pas \co{=} \\
\co{.first()} & un objet ou \co{None} \\
modifier puis \co{commit()} & \co{UPDATE}, pas besoin d'\co{add} \\
\co{db.delete(u)} + \co{commit()} & \co{DELETE} (\co{del u} : rien) \\
\bottomrule
\end{tabularx}

\begin{lstlisting}
user = db.query(User).filter(User.id == user_id).first()
if not user:
    raise HTTPException(status_code=404,
                        detail="Utilisateur non trouve")
\end{lstlisting}

\subsection{\co{response\_model}} valide \textbf{et filtre} la réponse : un
champ absent du schéma n'est pas envoyé. Dans le cours, \co{/users} avec
\co{UserSchema} ne renvoie donc \textbf{pas} l'\co{id}. Utile aussi pour ne
jamais exposer un mot de passe.

% =====================================================================
\section{Exercice : un-à-plusieurs}

\begin{lstlisting}
class User(Base):
    ...
    posts = relationship("Post", back_populates="author",
                         cascade="all, delete-orphan")

class Post(Base):
    __tablename__ = "posts"
    id = Column(Integer, primary_key=True, index=True)
    title = Column(String)
    content = Column(Text)
    user_id = Column(Integer, ForeignKey("users.id"),
                     nullable=False)
    author = relationship("User", back_populates="posts")
\end{lstlisting}

\begin{itemize}
  \item La \co{ForeignKey} va du côté \textbf{plusieurs} (le post).
  \item \co{ForeignKey("users.id")} : la \textbf{table} ;
        \co{relationship("Post")} : la \textbf{classe}.
  \item \co{relationship} n'ajoute \textbf{aucune colonne} : il donne
        \co{user.posts} et \co{post.author}.
  \item \co{cascade} : supprimer l'auteur supprime ses posts.
  \item SQLite n'applique les clés étrangères qu'avec
        \co{PRAGMA foreign\_keys=ON}.
  \item Routes : \co{POST /addpost}, \co{GET /posts},
        \co{GET /post/\{id\}}, \co{GET /user/\{id\}/posts},
        \co{PUT /post/\{id\}}, \co{DELETE /post/\{id\}}.
\end{itemize}

\end{multicols}

\vspace{0.3em}
\begin{astucebox}{Checklist des 60 dernières secondes}
\begin{minipage}[t]{0.485\linewidth}
\begin{itemize}
  \item \co{uvicorn main:app --reload}, port \textbf{8000}, \co{/docs}.
  \item Sans type : chaîne. Avec \co{int} invalide : \textbf{422}.
  \item Path = \co{\{\}}, query = type simple, body = Pydantic.
  \item \co{add} $\to$ \co{commit} $\to$ \co{refresh}.
\end{itemize}
\end{minipage}\hfill
\begin{minipage}[t]{0.485\linewidth}
\begin{itemize}
  \item \co{raise HTTPException(404)}, jamais \co{return}.
  \item \co{Depends(get\_db)} sans parenthèses ; \co{yield} + \co{finally}.
  \item \co{response\_model} filtre les champs.
  \item \co{ForeignKey} côté \og plusieurs \fg, nom de \textbf{table}.
\end{itemize}
\end{minipage}
\end{astucebox}

\vspace{0.4em}
\begin{astucebox}{Vocabulaire des énoncés (le QCM est en anglais)}
\begin{minipage}[t]{0.485\linewidth}
\begin{itemize}
  \item \emph{endpoint / route} : méthode + chemin
  \item \emph{path parameter} : partie variable de l'URL
  \item \emph{query parameter} : ce qui suit \co{?}
  \item \emph{request body} : le JSON envoyé
  \item \emph{dependency injection} : \co{Depends}
\end{itemize}
\end{minipage}\hfill
\begin{minipage}[t]{0.485\linewidth}
\begin{itemize}
  \item \emph{ORM} : Object-Relational Mapper
  \item \emph{session} : unité de travail SQLAlchemy
  \item \emph{schema} : modèle Pydantic
  \item \emph{one-to-many} : un-à-plusieurs
  \item \emph{foreign key} : clé étrangère
\end{itemize}
\end{minipage}
\end{astucebox}

\vspace{0.4em}
\begin{center}
\footnotesize\itshape
Format du contrôle : 6 propositions, 1 ou 2 bonnes réponses, une case fausse
annule la question.\\
Pour s'entraîner : \texttt{cd "Lesson 5" \&\& python quiz.py}
\end{center}

\end{document}
